\documentclass{article}
\usepackage{array}
\usepackage[a4paper, top=3cm, bottom=2.5cm, left=2.5cm, right=2.5cm]{geometry} % Ajuste de márgenes
\usepackage[spanish]{babel}
\usepackage[utf8]{inputenc}
\usepackage{tikz}
\usepackage{titling}
\usepackage{graphicx}
\usepackage{fancyhdr}
\usepackage{amsmath}
\usepackage{amssymb}
\usepackage{multicol}
\usepackage{cancel}
\usepackage{pgfplots}
\usepackage{hyperref}
\usepackage{bookmark}
\pgfplotsset{compat=1.18}
\usepackage{titlesec} % Para personalizar títulos
\usepackage{tocloft}  % Para mejorar el índice
\usepackage{setspace} % Para controlar el espaciado

\usepackage{xcolor}
\usepackage{enumitem}

\definecolor{headerblue}{RGB}{50,90,140}
\definecolor{responsegray}{RGB}{80,80,80}



% Configuración de Fancyhdr para encabezados y pies de página
\pagestyle{fancy}
\fancyhf{}
\fancyhead[L]{\includegraphics[width=2cm]{assets/logo-utp.png}}
\fancyhead[R]{\textbf{Diseño de Productos y Servicios}}

\fancyfoot[R]{\thepage} % Número de página alineado a la derecha

% Ajustes de espaciado entre párrafos y márgenes superiores
\setlength{\parskip}{1.5em}
\setlength{\parindent}{0pt}
\setlength{\headheight}{17.26935pt} % Altura del encabezado
\addtolength{\topmargin}{-2.26935pt} % Compensar el aumento de la altura del encabezado
\setlength{\textheight}{23cm}  % Ajusta el alto del texto

% Definición de comandos personalizados
\newcommand{\SubItem}[1]{
    {\setlength\itemindent{15pt} \item[-] #1}
}

% Título del documento con mejor control de espaciado
\title{
  \includegraphics[width=5cm]{./assets/logo-utp.png} \\
  \vspace{1cm}
  \textbf{Universidad Tecnológica del Perú} \\
  \vspace{2cm}
  \textbf{Avance de Proyecto 1: Sistema de Reservas de Salas de Estudio.} \\
  \vspace{1cm}
  \large \textbf{Para el curso de Diseño de Productos y Servicios.} \\
  \large \textbf{Ing. Software y Sistemas.}
}
\author{
  \textbf{Luis Huatay Salcedo.} \\
  \textbf{Garcia Chumpitaz Cindel Roxell.} \\
  \textbf{Díaz Benítez, Fernando Raúl.} \\
  \textbf{Quispe Fernandez, Bryan Alexander.}
}


\begin{document}
\maketitle
\begin{center}
  Sección 44698
\end{center}
\thispagestyle{empty}
\begin{center}
  Mg. Marcos Teodoro Yerren Huima  
\end{center}
\restoregeometry

\newpage
\begin{center}
  \textbf{\Large Índice}
\end{center}
\vspace{0.5cm} % Espacio entre título y contenid
\begin{spacing}{1.15} % Espaciado personalizado para mayor legibilidad
  \noindent
  \begin{enumerate}
    \item Introducción
    \item Definición del problema
    \item Objetivo
    \item Investigación de usuarios
  %   \begin{enumerate}
  %     \item Objetivos específicos
  %   \end{enumerate}
  %   \item Términos estadísticos
  %   \item Recolección de información
    \end{enumerate}
\end{spacing}

\newpage
\vspace*{\fill}
\section{Introducción}
La Universidad Tecnológica del Perú (UTP) es una institución educativa que ofrece una amplia gama de programas académicos y servicios a sus estudiantes. Uno de los aspectos fundamentales para el éxito académico es la disponibilidad de espacios adecuados para el estudio y la colaboración. En este contexto, el presente informe se centra en la creación de un sistema de reservas de salas de estudio, diseñado específicamente para satisfacer las necesidades de los estudiantes y mejorar su experiencia educativa.

El sistema de reservas de salas de estudio tiene como objetivo facilitar la gestión y el uso eficiente de los espacios de estudio disponibles en la UTP. A través de este sistema, los estudiantes podrán reservar salas de estudio de manera rápida y sencilla, asegurando así un entorno propicio para el aprendizaje y la colaboración.
\vspace*{\fill}

\newpage

\section{Definición del problema}

Muchos estudiantes tienen dificultades para reservar salas de estudio en la biblioteca 
debido a la falta de un sistema digital eficiente. Esto genera desorganización, pérdida de tiempo y conflictos por el uso de espacios.

La UTP cuenta con una biblioteca que ofrece diversas salas de estudio, pero la gestión de reservas que se realiza por su aplicación se hace de manera que una sala está completamente ocupada por el tiempo que se reserva y no por el tiempo que se ocupa, lo que puede llevar a confusiones y malentendidos entre los estudiantes. Además, la falta de un sistema centralizado dificulta la planificación y el uso óptimo de estos espacios.

La situación actual presenta varios problemas, entre ellos:
\begin{itemize}
    \item Falta de un sistema digital eficiente para la gestión de reservas.
    \item Desorganización en la utilización de los espacios de estudio.
    \item Conflictos entre estudiantes por el uso de las salas.
    \item Dificultad para encontrar y reservar salas de estudio disponibles.
    \item Pérdida de tiempo al buscar espacios adecuados para estudiar.
    \item Falta de información sobre la disponibilidad de salas en tiempo real.
\end{itemize}

\subsection{¿Qué problema o necesidad busca solucionar su proyecto?}

El proyecto busca implementar una plataforma web o app que permita a los estudiantes 
reservar salas de forma rápida, ordenada y sobre todo, que estas se mantengan disponibles a penas se desocupan. Esto no solo mejorará la experiencia de los estudiantes, sino que también optimizará el uso de los espacios de estudio en la UTP.

\newpage

\section{Investigación de usuarios}

La investigación de usuarios es un proceso fundamental para comprender las necesidades y expectativas de los estudiantes en relación con el sistema de reservas de salas de estudio. A través de encuestas y entrevistas, se recopilará información valiosa que guiará el diseño y desarrollo del sistema.

\subsection{Preguntas de entrevista}

Las preguntas de entrevista se caracterizan por ser abiertas y permitir a los usuarios expresar sus opiniones y experiencias de manera detallada. A continuación, se presentan algunas preguntas que se utilizarán en las entrevistas:

\begin{enumerate}
  \item ¿Cuánto tiempo sueles usar una sala de estudio por sesión?
  \item ¿Qué dificultades tienes al acceder a las salas de estudio? ¿Por qué?
  \item ¿Qué funciones te parecerían más útiles en un sistema de reservas para las salas de estudio?
  \item ¿Qué tan fácil te resulta usar el sistema de reservas y qué cambiarías para mejorarlo?
  \item ¿Crees que un sistema de reservas eficiente influye en tus actividades académicas? ¿Por qué?
\end{enumerate}

\subsection{Preguntas de encuesta}

Esta por otro lado, se caracteriza por ser cerrada y permite obtener datos cuantitativos. A continuación, se presentan algunas preguntas que se utilizarán en la encuesta:

\begin{multicols}{2}
  \begin{enumerate}
    \item ¿Tu universidad cuenta con un sistema de reservas para las salas de estudio? 
    \begin{itemize}
      \item $\square$ a) Sí
      \item $\square$ b) No
    \end{itemize}
  
    \item ¿Con qué frecuencia utilizas las salas de estudio en la universidad? 
    \begin{itemize}
      \item $\square$ a) Nunca
      \item $\square$ b) A veces
      \item $\square$ c) Frecuentemente
      \item $\square$ d) Siempre
    \end{itemize}
  
    \item ¿Sueles utilizar la sala de estudio de manera individual o en grupo? 
    \begin{itemize}
      \item $\square$ a) Individual
      \item $\square$ b) En grupo
      \item $\square$ c) Ambos
    \end{itemize}
  
    \item ¿En qué horarios puedes usar las salas de estudio? 
    \begin{itemize}
      \item $\square$ a) Mañana
      \item $\square$ b) Tarde
      \item $\square$ c) Noche
    \end{itemize}
  
    \item ¿Has tenido dificultades para encontrar una sala de estudio disponible? 
    \begin{itemize}
      \item $\square$ a) Sí
      \item $\square$ b) No
    \end{itemize}
  
    \item ¿Te gustaría contar con una aplicación o plataforma digital para reservar salas de estudio? 
    \begin{itemize}
      \item $\square$ a) Sí, definitivamente
      \item $\square$ b) Tal vez
      \item $\square$ c) No
    \end{itemize}
  \end{enumerate}  
\end{multicols}

\newpage
\section{Mapa de empatía}

Esta sección pretende resumir la información obtenida de los usuarios a través de las entrevistas y encuestas. El mapa de empatía es una herramienta visual que ayuda a comprender mejor a los usuarios y sus necesidades. A continuación, se presenta un ejemplo de un mapa de empatía para el sistema de reservas de salas de estudio:

\begin{center}
  \begin{tabular}{|>{\centering\arraybackslash}m{3cm}|>{\centering\arraybackslash}p{11cm}|}
    \hline
    \textbf{¿Qué piensa y siente?} & 
    \begin{itemize}[leftmargin=*, label={--}]
      \item Deseo de la implementación de un sistema de reservas de sala de estudio.
      \item Requiere estudiar en un ambiente tranquilo, sin bulla ni distracciones.
      \item No estar haciendo cola para pedir una reserva de manera presencial.
    \end{itemize} \\ \hline
    \textbf{¿Qué ve?} & 
    \begin{itemize}[leftmargin=*, label={--}]
      \item Sistemas que tienen poco diseño orientado al usuario.
      \item Dificultades para guiarse dentro de la plataforma.
      \item Solicitud de los estudiantes por mejoras en los sistemas de reserva.
    \end{itemize} \\ \hline
    \textbf{¿Qué escucha?} & 
    \begin{itemize}[leftmargin=*, label={--}]
      \item Comentarios de otros estudiantes con problemas similares.
      \item Sugerencias para añadir más funciones al sistema.
      \item Preferencias de sesiones de estudio grupales frente a las individuales.
    \end{itemize} \\ \hline
    \textbf{¿Qué dice y hace?} & 
    \begin{itemize}[leftmargin=*, label={--}]
      \item De manera general, el uso de las salas es ocasionalmente.
      \item Solicita mejoras en la interfaz del sistema.
      \item Preferencias por las salas de estudio en el turno de mañana.
    \end{itemize} \\ \hline
    \textbf{¿Cuáles son sus esfuerzos?} & 
    \begin{itemize}[leftmargin=*, label={--}]
      \item Dificultades para conseguir una sala de estudio por motivo de ocupación.
      \item Demora del sistema en conseguir salas de estudio.
      \item Problemas con la experiencia del usuario.
    \end{itemize} \\ \hline
    \textbf{¿Qué resultados espera?} & 
    \begin{itemize}[leftmargin=*, label={--}]
      \item Reservas garantizadas en todos los horarios.
      \item Mejoras en el sistema como notificaciones cuando el tiempo de uso de la sala culmine.
      \item Mejor gestión del uso de las salas.
    \end{itemize} \\ \hline
  \end{tabular}
\end{center}

\newpage

\section{Perfil de usuario y protopersona}

El perfil de usuario es una representación ficticia de un usuario típico del sistema de reservas de salas de estudio. A continuación, se presenta un perfil de usuario que resume las características y necesidades de los estudiantes, además de una protopersona que representa a un usuario específico:

\subsection{Perfil de Usuario}

\textbf{Sistema de Reservas de Salas en la Biblioteca:} Estudio universitario \\
\textbf{Edad:} 18 a 26 años \\
\textbf{Profesión:} Sin especificar \\
\textbf{Necesidad principal:} El estudiante requiere un espacio de estudio que garantice la concentración y el aprendizaje, libre de interrupciones externas, y que además le permita reservar dicha sala de manera remota sin necesidad de acudir a las oficinas del Servicio Estudiantil.

\subsection{Proto-persona}

\textbf{Nombre:} Camilo Soza \\
\textbf{Edad:} 20 años \\
\textbf{Carrera:} Derecho \\
\textbf{Motivaciones:} Le gusta la investigación de temas referentes al derecho penal. \\
\textbf{Frustraciones:} El estudiante se siente frustrado por la falta de un espacio de estudio adecuado dentro de la universidad para realizar sus actividades, y le resulta incómodo esperar largos periodos en la cola del Servicio Estudiantil, donde otros estudiantes realizan consultas diversas.

\end{document}